\section{Recolección de Datos}

\subsection{Plan de Recolección}

El proceso de recolección de datos se realizó de forma digital por medio de
la plataforma Google Forms, ésta última elegida debido a su simplicidad de uso
ya que es conocida por la mayoría de personas, haciendo que sea sencillo para
los encuestados el utilizarla. 

La encuesta se distribuyó por medios conocidos como redes sociales y
contactos personales conocidos por su simple distribución. 
Se pretendió que el público objetivo se alcanzó por los medios
más populares: Instagram, Facebook, WhatsApp. 

EL proceso fue sencillo: se envió la encuesta; el usuario es informado del fin de
la misma, además de que los datos registrados fueron anónimos y únicamente con
fines académicos; la encuesta no solicitó ningún dato personal justo porque no
es relevante para el proyecto; posterior a este proceso iniciaron las preguntas.
La encuesta fue dividida en seis secciones:
\begin{enumerate}
    \item Conocimientos sobre términos claves
    \item Comportamiento de compra
    \item Toma de decisiones
    \item Reacción ante estímulos ublicitarios
    \item Durabilidad
    \item Reparación
    \item Percepción de utilidad de la herramienta propuesta.
\end{enumerate}

Al finalizar la encuesta, se agradeció al encuestado el tiempo brindado y se le
reiteró el fin de recopilar estos datos, además de ofrecer la opción de
compartir el avance de la investigación a quien guste leerlo y, como cierre, una
reflexión para que no deje de cuestionarse, como consumidores, porqué compran lo
que compran y el gran poder de sus decisiones. 

\subsection{Plan de Procesamiento}
Al finalizar el trabajo de campo, los datos obtenidos pasaron por un proceso
sistemático de análisis orientado a verificar la hipótesis propuesta, el cual se
estructuró en dos etapas:

\begin{itemize}
    \item \textbf{Procesamiento cuantitativo}: Las respuestas provenientes de las encuestas se
          organizarán en una matriz de datos mediante software de hoja de cálculo (como
          Microsoft Excel). Antes del análisis, se llevó a cabo una depuración de la
          información para eliminar registros incompletos o que presenten inconsistencias.
          A continuación, se emplearon técnicas de estadística descriptiva con el fin de
          obtener frecuencias, porcentajes y medidas de tendencia central. 
          
    \item \textbf{Análisis de mercado}: Los datos organizados se tradujeron en representaciones
          visuales, tales como gráficas de barras y de pastel. Esto facilitó la
          interpretación del grado de aceptación hacia la plataforma, así como la
          relevancia que los encuestados atribuyeron a la información presentada, según lo
          que pudo inferirse de sus respuestas y comportamiento general. 
\end{itemize}

\subsection{Plan de Presentación}

Con el propósito de facilitar la interpretación y el análisis de la información
recolectada a lo largo de la investigación, se aplicó una encuesta a 241
personas, cuyos resultados fueron presentados mediante recursos
gráficos que permitieron visualizar de forma clara las tendencias identificadas. 

La presentación gráfica de los datos permitió contrastar los hallazgos
obtenidos con los objetivos e hipótesis establecidos, lo que facilitó la
elaboración de conclusiones y recomendaciones pertinentes.